\documentclass[aspectratio=169]{beamer}

\usetheme{Madrid}
\usecolortheme{default}

\usepackage[utf8]{inputenc}
\usepackage[T1]{fontenc}
\usepackage{amsmath,amssymb}
\usepackage{booktabs}

\title{Proof of Concept: Use and Fine-Tune MACE}
\date{\today}

\begin{document}

\begin{frame}
  \titlepage
\end{frame}

\begin{frame}{Outline}
  \tableofcontents
\end{frame}

\section{Dataset we already have}

\begin{frame}{Examples of configuration-group names (catalog)}
  \small
  \begin{columns}[T,onlytextwidth]
    \begin{column}{0.49\textwidth}
      \textbf{LiF / Li-rich}
      \begin{itemize}\setlength\itemsep{0.15em}
        \item \texttt{LiF64\_kjpaw}, \texttt{LiF64\_kjpaw\_v2}
        \item \texttt{LiF64\_kjpaw\_NPT\_v3}, \texttt{LiF64\_kjpaw\_final}
        \item \texttt{LiFinterface\_kjpaw\_v1}, \texttt{LiFinterface\_kjpaw\_v2}
        \item \texttt{LiwithF\_NPT\_final}, \texttt{LiwithF\_isolated}
        \item \texttt{LiwithBF}, \texttt{LiwithBF\_NPT\_final}
      \end{itemize}
    \end{column}
    \begin{column}{0.49\textwidth}
      \textbf{Salts, B--Li, cells, ionics}
      \begin{itemize}\setlength\itemsep{0.15em}
        \item \texttt{LiBF4}, \texttt{LiBF4\_NPT\_final}, \texttt{LiBF4v4}
        \item \texttt{LiBF4\_relax}
        \item \texttt{BLi3\_interface\_NPT}, \texttt{BLi3\_isolated}
        \item \texttt{BCC\_54\_kjpaw\_NPT}, \texttt{BCC54\_isolated}
        \item \texttt{EMIM\_TFSI\_BF4}, \texttt{EMIMLiTFSI\_anode\_rv2}
        \item \texttt{special}
      \end{itemize}
    \end{column}
  \end{columns}
\end{frame}

\section{FitSNAP model}

\section{FitSNAP model}
\begin{frame}{The FitSNAP model we have}
  \begin{itemize}
    \item \textbf{FitSNAP stack} in words:
      \begin{itemize}
        \item \textbf{Scraper} --- reads quantum-chemistry frames (e.g.\ JSON) into structures.
        \item \textbf{Calculator} --- builds \textbf{SNAP bispectrum} descriptors per atom in LAMMPS (radial cutoff, \texttt{twojmax}, chemical channels, etc.).
        \item \textbf{Solver} --- \textbf{PyTorch} MLP.
        \item \textbf{Layers:} PyTorch MLP \texttt{num\_desc $\to$ 256 $\to$ 128 $\to$ 64 $\to$ 64 $\to$ 1}, \texttt{multi\_element\_option = 2} (per-element nets).
    \item \textbf{Training:} learning rate \texttt{5e-6}, \texttt{70} epochs, batch \texttt{50}.
    \end{itemize}
  \end{itemize}
\end{frame}

\begin{frame}{Metrics}
    \begin{figure}
        \centering
        \includegraphics[width=0.8\linewidth]{Metrics MLPDFT.png}
        \caption{Metrics obtained by two models}
    \end{figure}
\end{frame}

\section{MACE model}
\begin{frame}{What the MACE model is}
    \begin{figure}
        \centering
        \includegraphics[width=0.8\linewidth]{image.png}
        \caption{I.~Batatia et al. 2022}
    \end{figure}
\end{frame}

\begin{frame}{Message Passing Neural Networks}
\begin{figure}
    \centering
    \includegraphics[width=0.8\linewidth]{mpnn.png}
    \caption{Message Passing Neural Network Concept}
\end{figure}
\end{frame}

\begin{frame}{MACE Step (1): Atom Embedding}
  A molecule or crystal is represented as a \textbf{graph}:
  atoms are nodes, edges connect pairs within a cutoff radius $r_c$.
  \medskip
  Each atom $i$ is assigned an initial feature vector:
  \[
    \mathbf{h}_i^{(0)} = \mathbf{W}_{Z_i}
  \]
  a learned embedding indexed by atomic number $Z_i$ (lookup table).
  \medskip
  \textbf{Goal:} produce, for each atom, a scalar energy $E_i$
  that reflects its local chemical environment.
  The total energy is $E = \sum_i E_i$.
\end{frame}

\begin{frame}{Depth on MPNN}

\begin{itemize}
    \item $t=0$, the atom/node only can see itself.
    \item $t=1$, the atom can see its neighborhoods - \textbf{TWO BODY}.
    \item $t=2$, the atom can see the neighborhoods from its neighborhoods - TRIPLETS - \textbf{THREE BODY}
\end{itemize}
Three and higher order bodies are computationally expensive.
\end{frame}
\begin{frame}{MACE Step (2): Interaction Block}
  For each atom $i$, messages from neighbors $j$ are built as:
  \[
    \mathbf{m}_{ij} = \sum_{k} W_k \, R_{k}(r_{ij}) \cdot
    Y_l^m\!\left(\hat{r}_{ij}\right) \otimes \mathbf{h}_j^{(t)}
  \]
  $W_k$ learned matrix.
  \begin{itemize}
    \item $Y_l^m(\hat{r}_{ij})$: \textbf{spherical harmonics} — fixed geometric
          functions encoding the \emph{direction} to neighbor $j$.
    \item $R_k(r_{ij})$: \textbf{radial network} — a small MLP acting on the
          scalar distance. \emph{This is where most parameters live.}
    \item $\otimes$: tensor product combining neighbor features with
          edge geometry, preserving equivariance via
          Clebsch--Gordan coefficients.
  \end{itemize}

  Messages are aggregated: $\mathbf{m}_i = \sum_{j \in \mathcal{N}(i)} \mathbf{m}_{ij}$
\end{frame}


\begin{frame}{MACE Step (3): Equivariant Product Block}
  After aggregating neighbors, atom $i$ holds a feature vector
  $\mathbf{m}_i$ encoding its local environment.

  \medskip
  \textbf{The problem:} one message passing layer only captures
  pairwise $(i,j)$ interactions --- that is \textbf{2-body}.

  \medskip
  \textbf{The MACE idea:} take tensor products of $\mathbf{m}_i$
  with itself $\nu$ times:
  \[
    \mathbf{m}_i \otimes \mathbf{m}_i \otimes \cdots
    \quad (\nu \text{ times})
  \]
  This mixes contributions from \emph{different} neighbors $j, k$
  simultaneously, producing \textbf{$(\nu+1)$-body correlations}
  without any extra message passing.

  \medskip
  \begin{itemize}
    \item With $\nu = 2$: reaches \textbf{3-body} after layer 1,
          compounding across layers.
    \item The CG coefficients ensure the result stays
          \textbf{equivariant} — no parameters, fixed by symmetry.
  \end{itemize}
\end{frame}


\begin{frame}{MACE Step (4): Readout}
  After $T$ interaction layers, each atom $i$ holds a feature vector
  $\mathbf{h}_i^{(t)}$ with components of different equivariance order $L$.
  \medskip
  Each layer contributes a partial atomic energy via its $L=0$ features:
  \[
    E_i = E_i^{(0)} + E_i^{(1)} + \cdots + E_i^{(T)}
  \]
  where each $E_i^{(t)}$ is a linear projection of $\mathbf{h}_i^{(t)}\big|_{L=0}$,
  except the last layer which uses a small MLP.

  \medskip
  The total energy and forces are then:
  \[
    E = \sum_i E_i, \qquad
    \mathbf{F}_i = -\frac{\partial E}{\partial \mathbf{r}_i}
  \]

  \textbf{Forces are free} --- no extra network needed, just
  automatic differentiation through the entire graph.
\end{frame}

\begin{frame}{Parameters Distribution}
    \begin{figure}
        \centering
        \includegraphics[width=0.7\linewidth]{image_par.png}
    \end{figure}
    \begin{figure}
        \centering
        \includegraphics[width=0.5\linewidth]{image_par2.png}
    \end{figure}
\end{frame}


\begin{frame}{MACE Flavors: Training Dataset}
MACE MP OMAT Number of parameters: 9063204
MACE MP 0 - Small - Number of parameters: 3847696
MACE MP 0B3 - Number of parameters: 9063204

Training data set:


\end{frame}

\section{Zero-shot evaluation}
\begin{frame}{Zero-Shot MACE-MP Evaluation on LiF 54}

\end{frame}


\begin{frame}{Challenges of Fine-Tuning MACE on a Specific Dataset}

  \begin{itemize}

    \item \textbf{Reference energy mismatch ($E_0$ problem)}\\
      The foundation model and the target DFT dataset use different
      atomic reference energies. A linear regression over the
      training set is required to align them before fine-tuning.

    \medskip
    \item \textbf{DFT functional mismatch}\\
      MACE-MP was trained with very specific data, if different
      functional (DFT), the model must \emph{unlearn} systematic
      biases --- which can hurt generalization.
    \medskip
    \item \textbf{Out-of-domain chemistry}\\
      LiBF$_4$ contains ionic species and Li, which are outside
      the MACE-OFF23 organic training domain. The model has no
      prior knowledge of ionic interactions, charge transfer,
      or metal coordination.
\end{itemize}
\end{frame}
\begin{frame}{Challenges}
    \begin{itemize}
    \item \textbf{Data scarcity}\\
      Fine-tuning on a small dataset risks \textbf{catastrophic
      forgetting} --- the model overfits to the new data and
      loses the general knowledge from pretraining.

    \item \textbf{Force/energy weight balance}\\
      The loss $\mathcal{L} = \lambda_E \mathcal{L}_E +
      \lambda_F \mathcal{L}_F$ is sensitive to the choice of
      $\lambda_E, \lambda_F$. Wrong balance leads to models
      that predict energies well but produce unphysical forces
      or vice versa.

  \end{itemize}


\end{frame}


\begin{frame}{Iterative Learning}
   \begin{figure}
        \centering
        \includegraphics[width=0.45\linewidth]{iterative_training.png}
    \end{figure}
\end{frame}

\begin{frame}{Active Learning}
\begin{figure}
    \centering
    \includegraphics[width=0.45\linewidth]{pic_044.png}
\end{figure}

\end{frame}

\end{document}
